\chapter{Formal definition}
\label{ch:formal}

This chapter defines Forge's syntax, its static semantics (which programs
are well typed) and its dynamic semantics (what they compute). The
definitions describe the core language that the front end produces: blocks
are already nested \fc|let|s, a function with several parameters is a chain
of one-parameter functions, operators are applications, and macros, modules
and imports have been expanded.

\section{Syntax}
\label{sec:grammar}
\index{grammar}

The grammar is written in EBNF: \nt{x} is a non-terminal, $[\,x\,]$ is
optional, $\{\,x\,\}$ repeats zero or more times, and terminals are set in
typewriter type.

\subsection*{Programs and items}
\begin{ebnf}
\prodn{program}{$\{$ \nt{item} $\}$}
\prodn{item}{\kw{let} \nt{binding}}
\alt{\kw{let rec} \nt{binding} $\{$ \kw{and} \nt{binding} $\}$}
\alt{\kw{type} $[$\kw{rec}$]$ \nt{typedecl} $\{$ \kw{and} \nt{typedecl} $\}$}
\alt{\nt{trait} $|$ \nt{impl}}
\alt{\kw{mod} \nt{Name} \kw{where} $\{$ \nt{item} $\}$ \kw{end}}
\alt{\kw{use} \nt{path} $|$ \kw{import} \nt{string}}
\alt{\kw{macro\_rules!} \nt{name} \kw{where} \nt{arm} $\{$ \tk{;} \nt{arm} $\}$ \kw{end}}
\prodn{binding}{\nt{pattern} $[$\tk{:} \nt{type}$]$ \tk{=} \nt{expr}}
\alt{\nt{name} $[$\tk{[} \nt{constraint} $\{$\tk{,} \nt{constraint}$\}$ \tk{]}$]$ $\{$ \nt{param} $\}$ $[$\tk{:} \nt{type}$]$ \tk{=} \nt{expr}}
\prodn{param}{\nt{subpattern} $|$ \tk{(} \nt{pattern} \tk{:} \nt{type} \tk{)}}
\prodn{typedecl}{\nt{Name} $[$\tk{<} \nt{tyvar} $\{$\tk{,} \nt{tyvar}$\}$ \tk{>}$]$ \tk{=} ( \nt{type} $|$ \nt{ctor} $\{$ \nt{ctor} $\}$ )}
\prodn{ctor}{\tk{|} \nt{Name} $[$ \kw{of} \nt{type} $]$}
\prodn{trait}{\kw{trait} \nt{Name} \tk{<} \nt{tparam} \tk{>} $[$\kw{requires} \nt{constraints}$]$ \kw{where} $\{$ \nt{titem} $\}$ \kw{end}}
\alt{\kw{inter} \nt{Name} \tk{<} \nt{tparam} \tk{>} \tk{\{} $\{$ \nt{titem} $\}$ \tk{\}}}
\prodn{tparam}{\nt{tyvar} $|$ \nt{tyvar} \tk{<\_>}}
\prodn{titem}{\kw{val} \nt{name} \tk{:} \nt{type} $|$ \kw{let} \nt{name} $\{$ \nt{param} $\}$ \tk{=} \nt{expr}}
\prodn{impl}{\kw{impl} \nt{Name} \kw{for} \nt{type} $[$\kw{requires} \nt{constraints}$]$ \kw{where} \nt{method} $\{$ \tk{,} \nt{method} $\}$ \kw{end}}
\prodn{method}{\nt{name} $\{$ \nt{param} $\}$ \tk{=} \nt{expr}}
\prodn{constraints}{\nt{Name}\tk{<}\nt{type}\tk{>} $\{$ \tk{,} \nt{Name}\tk{<}\nt{type}\tk{>} $\}$}
\end{ebnf}

\subsection*{Expressions}
\begin{ebnf}
\prodn{expr}{\kw{fn} \nt{param} $[$\tk{:} \nt{type}$]$ \tk{->} \nt{expr}}
\alt{\kw{if} \nt{expr} \kw{then} \nt{expr} \kw{else} \nt{expr}}
\alt{\kw{let} \nt{binding} \kw{in} \nt{expr} $|$ \kw{let rec} \nt{binding} $\{$\kw{and} \nt{binding}$\}$ \kw{in} \nt{expr}}
\alt{\kw{case} \nt{expr} \kw{do} \tk{|} \nt{pattern} \tk{->} \nt{expr} $\{$ \tk{|} \nt{pattern} \tk{->} \nt{expr} $\}$}
\alt{\nt{pipe}}
\prodn{pipe}{\nt{cons} $\{$ \tk{|>} \nt{cons} $\}$}
\prodn{cons}{\nt{or} $[$ \tk{::} \nt{cons} $]$}
\prodn{or}{\nt{and} $[$ \tk{||} \nt{or} $]$}
\prodn{and}{\nt{rel} $[$ \tk{\&\&} \nt{and} $]$}
\prodn{rel}{\nt{rel} \nt{relop} \nt{add} $|$ \nt{add}}
\prodn{add}{\nt{add} \nt{addop} \nt{mul} $|$ \nt{mul}}
\prodn{mul}{\nt{mul} \nt{mulop} \nt{app} $|$ \nt{app}}
\prodn{app}{\nt{app} \nt{factor} $|$ \nt{factor}}
\prodn{factor}{\nt{literal} $|$ \nt{name} $|$ \nt{path} $|$ \tk{-} \nt{factor} $|$ \nt{factor} \tk{.} \nt{name}}
\alt{\tk{(} \nt{expr} \tk{)} $|$ \tk{(} \nt{expr} \tk{,} \nt{expr} $\{$\tk{,} \nt{expr}$\}$ \tk{)} $|$ \tk{(} \nt{operator} \tk{)}}
\alt{\tk{[} $[$\nt{expr} $\{$\tk{,} \nt{expr}$\}]$ \tk{]} $|$ \tk{[} \nt{expr} \tk{...} \nt{expr} \tk{]}}
\alt{\tk{[} \nt{expr} \tk{|} \nt{gen} $\{$\tk{,} \nt{gen}$\}$ \tk{]}}
\alt{\tk{\{} \nt{name} \tk{:} \nt{expr} $\{$\tk{,} \ldots$\}$ \tk{\}} $|$ \tk{\{} \nt{expr} \kw{with} \nt{name} \tk{=} \nt{expr} $\{$\tk{,} \ldots$\}$ \tk{\}}}
\alt{\tk{\{} \nt{stmt} $\{$ \tk{;} \nt{stmt} $\}$ \tk{\}} $|$ \nt{name}\tk{!(} \nt{tokens} \tk{)} $|$ \nt{name}\tk{![} \nt{tokens} \tk{]} $|$ \nt{name}\tk{!\{} \nt{tokens} \tk{\}}}
\prodn{gen}{\nt{pattern} \tk{<-} \nt{expr} $|$ \nt{expr}}
\prodn{stmt}{\kw{let} \nt{binding} $|$ \kw{let rec} \nt{binding} $|$ \nt{expr}}
\end{ebnf}
\noindent where \nt{relop} is \tk{== != <> < > <= >=} or a custom operator
starting with \tk{< > = !}, \nt{addop} is \tk{+ - \^} or a custom operator
starting with one of them, and \nt{mulop} is \tk{* / \%} or a custom operator
starting with one of them.

\subsection*{Patterns and types}
\begin{ebnf}
\prodn{pattern}{\nt{subpattern} $[$ \tk{::} \nt{pattern} $]$}
\prodn{subpattern}{\nt{name} $|$ \tk{\_} $|$ \nt{literal} $|$ \tk{-}\nt{int} $|$ \tk{()} $|$ \tk{[]}}
\alt{\tk{(} \nt{pattern} $\{$\tk{,} \nt{pattern}$\}$ \tk{)} $|$ \tk{\{} \nt{name} $[$\tk{:} \nt{pattern}$]$ $\{$\tk{,} \ldots$\}$ \tk{\}}}
\alt{\nt{Name} $[$ \nt{subpattern} $]$}
\prodn{type}{\nt{tfactor} $[$ \tk{->} \nt{type} $]$}
\prodn{tfactor}{\tk{Int} $|$ \tk{Float} $|$ \tk{Bool} $|$ \tk{Char} $|$ \tk{String} $|$ \tk{Unit} $|$ \nt{tyvar}}
\alt{\nt{Name} $[$\tk{<} \nt{type} $\{$\tk{,} \nt{type}$\}$ \tk{>}$]$ $|$ \tk{[} \nt{type} \tk{]} $|$ \tk{(} \nt{type} $\{$\tk{,} \nt{type}$\}$ \tk{)}}
\alt{\tk{\{} \nt{name} \tk{:} \nt{type} $\{$\tk{,} \ldots$\}$ \tk{\}}}
\end{ebnf}

\section{Static semantics}
\label{sec:static-semantics}
\index{typing rules}

\paragraph{Types and schemes} Monotypes $\tau$ and type schemes $\sigma$ are
\[
\begin{array}{rcl}
\tau & ::= & \alpha \mid \ty{Int} \mid \ty{Float} \mid \ty{Bool} \mid \ty{Char} \mid \ty{String} \mid \ty{Unit}
  \mid \tau_1 \to \tau_2 \mid \ty{List}\,\tau \mid \tau_1 \times\dots\times \tau_n \\
 & \mid & T\,\tau_1\dots\tau_n \mid \{\ell_1:\tau_1,\dots,\ell_n:\tau_n \mid \rho\} \mid \{\ell_1:\tau_1,\dots,\ell_n:\tau_n\} \\
\sigma & ::= & \forall \alpha_1\dots\alpha_n.\; P \Rightarrow \tau \qquad\qquad P ::= C_1\,\tau_1, \dots, C_k\,\tau_k
\end{array}
\]
where $\rho$ is a \emph{row variable} standing for the remaining fields of
an open record, and $P$ is a set of trait predicates. An environment
$\Gamma$ maps names to schemes; the set $\mathcal{I}$ of instances maps
each $C\,\tau$ that holds to the instance that proves it. $\Gamma \vdash e :
\tau$ reads ``$e$ has type $\tau$ in $\Gamma$'', and
$\Gamma \vdash p : \tau \Rightarrow \Gamma'$ ``pattern $p$ matches values of
type $\tau$, binding $\Gamma'$''. $\mathrm{gen}_\Gamma(\tau)$ quantifies
the type variables of $\tau$ (and of its predicates) that are not free in
$\Gamma$; $\tau \preceq \sigma$ says that $\tau$ is an instance of $\sigma$
whose predicates are all in $\mathcal{I}$ or among the predicates assumed
for the current definition.

\paragraph{Expressions}
\ruleblock{
\infrule{Var}{x : \sigma \in \Gamma \qquad \tau \preceq \sigma}{\Gamma \vdash x : \tau}
\qquad
\infrule{Lit}{}{\Gamma \vdash 42 : \ty{Int} \quad \Gamma \vdash \texttt{"s"} : \ty{String} \quad \dots}
\\[14pt]
\infrule{Fn}{\Gamma \vdash p : \tau_1 \Rightarrow \Gamma' \qquad \Gamma, \Gamma' \vdash e : \tau_2}{\Gamma \vdash \kw{fn}\ p \to e : \tau_1 \to \tau_2}
\qquad
\infrule{App}{\Gamma \vdash e_1 : \tau_1 \to \tau_2 \qquad \Gamma \vdash e_2 : \tau_1}{\Gamma \vdash e_1\ e_2 : \tau_2}
\\[14pt]
\infrule{Let}{\Gamma \vdash e_1 : \tau_1 \qquad \Gamma \vdash p : \tau_1 \Rightarrow \{x_i : \tau_i\} \qquad \Gamma, \{x_i : \mathrm{gen}_\Gamma(\tau_i)\} \vdash e_2 : \tau}{\Gamma \vdash \kw{let}\ p = e_1\ \kw{in}\ e_2 : \tau}
\\[14pt]
\infrule{LetRec}{\Gamma, \{f_i : \tau_i\} \vdash e_i : \tau_i \quad (1 \le i \le n) \qquad \Gamma, \{f_i : \mathrm{gen}_\Gamma(\tau_i)\} \vdash e : \tau}{\Gamma \vdash \kw{let rec}\ f_1 = e_1\ \kw{and}\ \dots\ \kw{and}\ f_n = e_n\ \kw{in}\ e : \tau}
\\[14pt]
\infrule{If}{\Gamma \vdash e_1 : \ty{Bool} \qquad \Gamma \vdash e_2 : \tau \qquad \Gamma \vdash e_3 : \tau}{\Gamma \vdash \kw{if}\ e_1\ \kw{then}\ e_2\ \kw{else}\ e_3 : \tau}
\qquad
\infrule{Case}{\Gamma \vdash e : \tau \qquad \Gamma \vdash p_i : \tau \Rightarrow \Gamma_i \qquad \Gamma, \Gamma_i \vdash e_i : \tau'}{\Gamma \vdash \kw{case}\ e\ \kw{do}\ \mid p_i \to e_i : \tau'}
\\[14pt]
\infrule{Arith}{\Gamma \vdash e_1 : \ty{Int} \qquad \Gamma \vdash e_2 : \ty{Int}}{\Gamma \vdash e_1 \oplus e_2 : \ty{Int}}
\quad
\infrule{Compare}{\Gamma \vdash e_1 : \tau \qquad \Gamma \vdash e_2 : \tau}{\Gamma \vdash e_1 \lessgtr e_2 : \ty{Bool}}
\quad
\infrule{Concat}{\Gamma \vdash e_i : \ty{String}}{\Gamma \vdash e_1 \mathbin{\hat{}} e_2 : \ty{String}}
\\[14pt]
\infrule{Nil}{}{\Gamma \vdash [\,] : \ty{List}\,\tau}
\quad
\infrule{Cons}{\Gamma \vdash e_1 : \tau \qquad \Gamma \vdash e_2 : \ty{List}\,\tau}{\Gamma \vdash e_1 :: e_2 : \ty{List}\,\tau}
\quad
\infrule{Range}{\Gamma \vdash e_i : \ty{Int}}{\Gamma \vdash [e_1 \ldots e_2] : \ty{List}\,\ty{Int}}
\\[14pt]
\infrule{Comp}{\Gamma \vdash e_1 : \ty{List}\,\tau_1 \qquad \Gamma \vdash p : \tau_1 \Rightarrow \Gamma' \qquad \Gamma, \Gamma' \vdash [e \mid \overline{g}] : \ty{List}\,\tau}{\Gamma \vdash [e \mid p \leftarrow e_1, \overline{g}] : \ty{List}\,\tau}
\quad
\infrule{Guard}{\Gamma \vdash c : \ty{Bool} \qquad \Gamma \vdash [e \mid \overline{g}] : \ty{List}\,\tau}{\Gamma \vdash [e \mid c, \overline{g}] : \ty{List}\,\tau}
\\[14pt]
\infrule{Tuple}{\Gamma \vdash e_i : \tau_i}{\Gamma \vdash (e_1, \dots, e_n) : \tau_1 \times \dots \times \tau_n}
\quad
\infrule{Record}{\Gamma \vdash e_i : \tau_i}{\Gamma \vdash \{\ell_i : e_i\} : \{\ell_i : \tau_i\}}
\\[14pt]
\infrule{Field}{\Gamma \vdash e : \{\ell : \tau \mid \rho\}}{\Gamma \vdash e.\ell : \tau}
\quad
\infrule{Update}{\Gamma \vdash e : \{\ell_i : \tau_i \mid \rho\} \qquad \Gamma \vdash e_i : \tau_i}{\Gamma \vdash \{e\ \kw{with}\ \ell_i = e_i\} : \{\ell_i : \tau_i \mid \rho\}}
}

\noindent A closed record type $\{\overline{\ell:\tau}\}$ unifies with an
open one $\{\overline{\ell':\tau'} \mid \rho\}$ when every $\ell'$ is among
the $\ell$, and then $\rho$ stands for the remaining fields; two open record
types share a fresh row. \textsc{Compare} applies to every type except
function types: comparing functions is a run-time error.

\paragraph{Patterns}
\ruleblock{
\infrule{PVar}{}{\Gamma \vdash x : \tau \Rightarrow \{x : \tau\}}
\quad
\infrule{PWild}{}{\Gamma \vdash \_ : \tau \Rightarrow \emptyset}
\quad
\infrule{PLit}{\Gamma \vdash \ell : \tau}{\Gamma \vdash \ell : \tau \Rightarrow \emptyset}
\\[14pt]
\infrule{PCons}{\Gamma \vdash p_1 : \tau \Rightarrow \Gamma_1 \qquad \Gamma \vdash p_2 : \ty{List}\,\tau \Rightarrow \Gamma_2}{\Gamma \vdash p_1 :: p_2 : \ty{List}\,\tau \Rightarrow \Gamma_1 \uplus \Gamma_2}
\quad
\infrule{PTuple}{\Gamma \vdash p_i : \tau_i \Rightarrow \Gamma_i}{\Gamma \vdash (p_1,\dots,p_n) : \tau_1 \times\dots\times \tau_n \Rightarrow \biguplus \Gamma_i}
\\[14pt]
\infrule{PCtor}{C : \forall \overline{\alpha}.\, \tau' \to T\,\overline{\alpha} \in \Gamma \qquad \Gamma \vdash p : \tau'[\overline{\tau}/\overline{\alpha}] \Rightarrow \Gamma'}{\Gamma \vdash C\ p : T\,\overline{\tau} \Rightarrow \Gamma'}
\quad
\infrule{PRecord}{\Gamma \vdash p_i : \tau_i \Rightarrow \Gamma_i}{\Gamma \vdash \{\ell_i : p_i\} : \{\ell_i : \tau_i \mid \rho\} \Rightarrow \biguplus \Gamma_i}
}
\noindent $\uplus$ is disjoint union: a variable may be bound only once in a
pattern.

\paragraph{Traits} A trait declaration $\kw{trait}\ C\langle\alpha\rangle\
\kw{where}\ \kw{val}\ m : \tau\ \kw{end}$ adds $m : \forall \alpha\,
\overline{\beta}.\; C\,\alpha \Rightarrow \tau$ to $\Gamma$, where
$\overline{\beta}$ are the other variables of $\tau$. An instance
$\kw{impl}\ C\ \kw{for}\ \tau_0\ \kw{requires}\ P\ \kw{where}\ \dots$ is
checked by typing each method's equation at $\tau[\tau_0/\alpha]$ assuming $P$,
and adds the rule ``$P$ implies $C\,\tau_0$'' to $\mathcal{I}$. A top-level
definition whose type keeps predicates on its own type variables is
generalised with them; a predicate on a type without variables must follow
from $\mathcal{I}$, or the program is rejected (\textit{No instance}).

\section{Dynamic semantics}
\label{sec:dynamic-semantics}
\index{evaluation rules}

Evaluation is strict and goes from left to right. Values are
\[
\begin{aligned}
v ::={}& n \mid f \mid c \mid s \mid \kw{true} \mid \kw{false} \mid () \mid [\,] \mid v_1 :: v_2 \mid (v_1,\dots,v_n) \mid \{\ell_i = v_i\} \\
  \mid{}& C \mid C\,v \mid \langle \rho, p, e \rangle \mid \langle \rho, f, p, e \rangle^{\mathrm{rec}}
\end{aligned}
\]
the last two being closures and recursive closures. $\rho \vdash e \Downarrow v$
reads ``in environment $\rho$, $e$ evaluates to $v$'';
$\mathrm{match}(p, v) = \rho'$ is the environment that pattern $p$ binds when it
matches $v$, and is undefined otherwise.

\ruleblock{
\infrule{EVar}{\rho(x) = v}{\rho \vdash x \Downarrow v}
\quad
\infrule{EFn}{}{\rho \vdash \kw{fn}\ p \to e \Downarrow \langle \rho, p, e \rangle}
\\[14pt]
\infrule{EApp}{\rho \vdash e_1 \Downarrow \langle \rho', p, e \rangle \qquad \rho \vdash e_2 \Downarrow v_2 \qquad \rho', \mathrm{match}(p, v_2) \vdash e \Downarrow v}{\rho \vdash e_1\ e_2 \Downarrow v}
\\[14pt]
\infrule{EAppRec}{\rho \vdash e_1 \Downarrow c = \langle \rho', f, p, e \rangle^{\mathrm{rec}} \qquad \rho \vdash e_2 \Downarrow v_2 \\ \rho', f \mapsto c, \mathrm{match}(p, v_2) \vdash e \Downarrow v}{\rho \vdash e_1\ e_2 \Downarrow v}
\\[14pt]
\infrule{ELet}{\rho \vdash e_1 \Downarrow v_1 \qquad \rho, \mathrm{match}(p, v_1) \vdash e_2 \Downarrow v}{\rho \vdash \kw{let}\ p = e_1\ \kw{in}\ e_2 \Downarrow v}
\quad
\infrule{ELetRec}{\rho, f \mapsto \langle \rho, f, p, e_1 \rangle^{\mathrm{rec}} \vdash e_2 \Downarrow v}{\rho \vdash \kw{let rec}\ f = \kw{fn}\ p \to e_1\ \kw{in}\ e_2 \Downarrow v}
\\[14pt]
\infrule{EIfT}{\rho \vdash e_1 \Downarrow \kw{true} \qquad \rho \vdash e_2 \Downarrow v}{\rho \vdash \kw{if}\ e_1\ \kw{then}\ e_2\ \kw{else}\ e_3 \Downarrow v}
\quad
\infrule{EIfF}{\rho \vdash e_1 \Downarrow \kw{false} \qquad \rho \vdash e_3 \Downarrow v}{\rho \vdash \kw{if}\ e_1\ \kw{then}\ e_2\ \kw{else}\ e_3 \Downarrow v}
\\[14pt]
\infrule{ECase}{\rho \vdash e \Downarrow v \qquad \mathrm{match}(p_j, v) = \rho_j \qquad \forall i < j.\ \mathrm{match}(p_i, v)\ \text{undefined} \\ \rho, \rho_j \vdash e_j \Downarrow v'}{\rho \vdash \kw{case}\ e\ \kw{do}\ \mid p_1 \to e_1 \dots \mid p_n \to e_n \Downarrow v'}
\\[14pt]
\infrule{EAndF}{\rho \vdash e_1 \Downarrow \kw{false}}{\rho \vdash e_1 \mathbin{\&\&} e_2 \Downarrow \kw{false}}
\quad
\infrule{EAndT}{\rho \vdash e_1 \Downarrow \kw{true} \qquad \rho \vdash e_2 \Downarrow v}{\rho \vdash e_1 \mathbin{\&\&} e_2 \Downarrow v}
\\[14pt]
\infrule{EDiv}{\rho \vdash e_1 \Downarrow n_1 \qquad \rho \vdash e_2 \Downarrow n_2 \qquad n_2 \ne 0}{\rho \vdash e_1 / e_2 \Downarrow \mathrm{trunc}(n_1 / n_2)}
\\[14pt]
\infrule{ERange}{\rho \vdash e_1 \Downarrow a \qquad \rho \vdash e_2 \Downarrow b}{\rho \vdash [e_1 \ldots e_2] \Downarrow [a, a+1, \dots, b]}
\quad
\infrule{ECmp}{\rho \vdash e_1 \Downarrow v_1 \qquad \rho \vdash e_2 \Downarrow v_2}{\rho \vdash e_1 < e_2 \Downarrow \mathrm{compare}(v_1, v_2) < 0}
\\[14pt]
\infrule{EComp}{\rho \vdash e_1 \Downarrow [v_1, \dots, v_k] \\ \rho, \mathrm{match}(p, v_i) \vdash [e \mid \overline{g}] \Downarrow \ell_i \ \text{for each $v_i$ that $p$ matches}}{\rho \vdash [e \mid p \leftarrow e_1, \overline{g}] \Downarrow \ell_{i_1} \mathbin{++} \dots \mathbin{++} \ell_{i_m}}
}

\noindent $\mathrm{compare}$ is the structural order of
Section~\ref{sec:operators}: by value for numbers and characters,
lexicographic for strings, lists and tuples, field by field (in name order)
for records, by constructor declaration order and then payload for sum types,
and undefined (a run-time error) for functions. A guard $c$ in a comprehension
keeps the current combination when $c \Downarrow \kw{true}$. When no rule
applies --- division by zero, a failed \fc|case| or refutable \fc|let|, a
comparison of functions --- evaluation stops with a run-time error.

A program evaluates its items in order: \fc|let p = e| evaluates $e$ in the
environment of the previous items and extends it with $\mathrm{match}(p, v)$.
Trait methods evaluate to the implementation of the instance chosen by the
types (\S\ref{sec:impl-traits}); the semantics above apply to the program
after that choice.
